\chapter{Role-graded Hessian и гейт embedding-нормировки}

Минимальная возможная лазейка в no-go версии II состоит в том, что
ненормированная фоновая деформация и нормированное квантовое состояние не
обязаны задаваться одним и тем же вектором конфигурационного пространства,
даже если измеряются одной Hodge-метрикой.

\section{Единая метрика}

Пусть нуль-формы всех блоков имеют общий вес \(w_0\), а интегральная
одна форма на \(\gamma\) --- вес \(w_1\). Для role-graded embeddings
фиксируем
\begin{align}
 \Xi_\tau
 &=1_{\mathbb{RP}^3}\oplus1_{S^1}\oplus P_{\perp n},\\
 \Xi_\nu
 &=P_H\otimes\widehat 1_{S^1}
   \oplus P_{\ker}\otimes e_1,
 \qquad \oint_\gamma e_1=1.
\end{align}
Здесь \(\widehat1_{S^1}=(2\pi)^{-1/2}\) является нормированным состоянием,
тогда как \(1_{S^1}\) в \(\Xi_\tau\) читается как единичная фоновая
амплитуда.

Одна factor-blind Hodge-метрика даёт
\begin{align}
 \|\Xi_\tau\|^2
 &=w_0\left(\pi^2+2\pi+\frac23\right),\label{eq:V3TauHessian}\\
 \|\Xi_\nu\|^2
 &=23w_0+\frac{w_1}{\pi}.\label{eq:V3NuHessian}
\end{align}
При standard Hodge choice
\begin{equation}
 w_0=w_1=1
\end{equation}
обе требуемые числовые структуры возникают одновременно. Следовательно,
чисто алгебраическая несовместимость отсутствует: no-go версии II зависит
от предположения, что обе коллективные координаты сначала переводятся в
\(L^2\)-нормированные состояния.

\begin{proposition}[Алгебраический двухсекторный pass]
Role-graded embedding и единая Hodge-метрика с единичными form-degree
весами одновременно дают
\begin{equation}
 \|\Xi_\tau\|^2=\pi^2+2\pi+\frac23,
 \qquad
 \|\Xi_\nu\|^2=23+\pi^{-1}.
\end{equation}
Независимые sector weights для этого тождества не требуются.
\end{proposition}

\section{Field-redefinition red-team}

Положительный числовой результат ещё не является физическим. Масштаб
интегральной формы фиксирован условием периода, а нормированное состояние
фиксировано вероятностной нормой. Но обычная вещественная нуль-форма
допускает замену координаты
\begin{equation}
 \phi_\tau\longmapsto a\phi_\tau.
\end{equation}
При этом Hessian-вклад меняется как \(a^{-2}\), а interaction vertex должен
измениться как соответствующая степень \(a\). Поэтому утверждение
``единичная фоновая амплитуда'' не является инвариантным, пока единица не
выведена из compact field range, charge lattice, canonical symplectic form
или нормированной vertex map.

Если разрешить embedding-scale \(a_\tau\), формула
\eqref{eq:V3TauHessian} принимает вид
\begin{equation}
 \|\Xi_\tau(a_\tau)\|^2
 =a_\tau^2w_0\left(\pi^2+2\pi+\frac23\right).
\end{equation}
Ни isometry group носителя, ни степень формы не фиксируют \(a_\tau=1\).

\begin{nogocriterion}[Embedding-normalization gap]
Role-graded Hessian проходит числовой двухсекторный тест, но не проходит
parent-action gate, пока масштаб фоновой charged-lepton coordinate не
получен из симметрии или канонической скобки и не показано совместное
преобразование её Hessian и vertex. Назначение \(a_\tau=1\) после просмотра
сырого seed является скрытой нормировкой.
\end{nogocriterion}

\section{Следующий конструктивный класс}

Следующий этап больше не ищет новую таблицу весов. Требуется классифицировать
механизмы, способные фиксировать \(a_\tau\):
\begin{enumerate}
 \item compact phase или holonomy coordinate с заранее заданным периодом;
 \item collective coordinate defect-а с нормировкой из symplectic form;
 \item boundary/AKSZ variable с квантованной pairing form;
 \item finite spectral triple, где масштаб задаётся representation trace и
 first-order constraints.
\end{enumerate}
Первым будет проверен compact-phase вариант, поскольку он минимально меняет
carrier и имеет немедленный отрицательный критерий: если требуемые объёмные
вклады превращаются в inverse-length stiffness или вводят decay constant,
маршрут закрывается.

Машинный аудит сохранён в
\texttt{s2t\_v3\_role\_graded\_hessian\_results.json}.